% Generated draft: requires mathematical and source review.
\documentclass[11pt]{article}
\usepackage[margin=25mm]{geometry}
\usepackage{fontspec}
\setmainfont{Latin Modern Roman}
\usepackage{amsmath,amssymb,mathtools,unicode-math}
\setmathfont{Latin Modern Math}
\usepackage{enumitem,needspace,xurl,hyperref}
\setlength{\parindent}{0pt}
\setlength{\parskip}{5pt}
\setlength{\emergencystretch}{3em}
\newcommand{\N}{\mathbb{N}}
\newcommand{\Z}{\mathbb{Z}}
\newcommand{\Q}{\mathbb{Q}}
\newcommand{\R}{\mathbb{R}}
\newcommand{\C}{\mathbb{C}}
\newcommand{\F}{\mathbb{F}}
\newcommand{\A}{\mathbb{A}}
\newcommand{\PP}{\mathbb{P}}
\newcommand{\E}{\mathbb{E}}
\newcommand{\eps}{\varepsilon}
\newcommand{\dd}{\,\mathrm{d}}
\begin{document}
\section*{UnsolvedMath ID: 30001537}
\textbf{Source problem code:} catalogue entry
\subsection{Definitions and mathematical statement}
A pro-$p$ group, for a prime $p$, is a topological group isomorphic to an inverse limit of finite $p$-groups. For an integer $g\geq 0$, a pro-$p$ presentation with $g$ generators and $g$ relators has the form
\[
G\cong F_g/\overline{\langle\!\langle r_1,\ldots,r_g\rangle\!\rangle},
\]
where $F_g$ is the free pro-$p$ group on $g$ generators, each $r_i$ lies in $F_g$, and the denominator is the closed normal subgroup generated by the relators. The source does not specify whether the generator and relator counts must be minimal, which values of $g$ are intended, or whether $p$ and $g$ are fixed as $S$ varies. Isomorphisms of pro-$p$ groups are understood here as continuous isomorphisms, and $\operatorname{Aut}(G)$ denotes continuous automorphisms; any intended restriction on these automorphisms remains unspecified.

The supplied question asks whether, for a given $g$-generator, $g$-relator pro-$p$ group $G$, varying a finite set $S$ of primes gives
\[
\Pr\bigl(G_S\cong G\bigr)\stackrel{?}{=}\frac{A(G)}{|\operatorname{Aut}(G)|}.
\]
This is an incomplete distributional formulation: the supplied record does not define $G_S$, the function $A(G)$, the admissible sets $S$, or the probability measure or limiting-density procedure represented by $\Pr$. Although the title refers to Galois groups, no base field, extension, or ramification conditions defining $G_S$ are supplied. Furthermore, a pro-$p$ group need not have a finite automorphism group, so the numerical meaning of the denominator and of the proposed ratio also requires clarification. These missing data prevent a fully specified conjecture; the displayed equality preserves the source question without supplying additional arithmetic hypotheses or a random-group model.
\subsection{Short English statement}
The question asks whether varying a finite set of primes produces a specified pro-$p$ group $G$ with probability $A(G)/|\operatorname{Aut}(G)|$. The supplied record leaves the arithmetic group, weighting function, sampling procedure, and interpretation of the ratio undefined, so this remains a statement-only draft requiring clarification and human mathematical review.
\subsection{Sources}
\begin{itemize}
\item Local UnsolvedMath registry, \texttt{data/unsolvedmath/problems.json}, permanent ID $30001537$, title ``Distribution of Profinite Galois Groups,'' source problem code \texttt{catalogue entry}. The supplied record belongs to the Oberwolfach Reports Open Problems collection. No individual proposer is identified; the recorded proposed year is 2010.
\item Supplied report citation: \emph{Low-Dimensional Topology and Number Theory} (2011), Oberwolfach Reports, pp.~2101--2163, DOI \texttt{10.4171/owr/2010/35}, \url{https://doi.org/10.4171/owr/2010/35}. The supplied background describes Oberwolfach Reports as documenting mathematical workshops and their open-problem discussions. The citation year 2011 and the record's 2010 year attribution are retained without reconciliation.
\item Supplied literature reference, labeled ``Boston--Ellenberg primary paper'': \url{https://doi.org/10.4171/GGD/127}. No further bibliographic details are supplied in the record.
\item Supplied literature reference, labeled ``original OWR report'': \url{https://doi.org/10.4171/owr/2010/35}.
\item Literature assessment checked on \texttt{2026-08-21}: this is supplied dated curation, not independently verified. The record labels the problem ``partially solved'' and states that the Boston--Ellenberg random pro-$p$ prediction has proofs and evidence in several cases and yields density results for associated quotients, while the general distribution conjecture in the extracted statement is not proved. The supplied background describes this assessment as dated literature triage derived from the OWR curation and requests independent verification of time-sensitive status and cited claims.
\item The supplied statement status is \texttt{corrected\_verified}; its accompanying verification describes only enclosing TeX notation in mathematical delimiters without changing mathematical content. This label does not establish a complete formulation or mathematical verification.
\end{itemize}
\end{document}
